\section{Bilangan Prima dan Pengujian Primalitas}

Bilangan prima adalah batu bata penyusun seluruh sistem kriptografi modern. Keamanan algoritma seperti RSA tidak didasarkan pada kerahasiaan algoritma, tetapi pada fakta matematis bahwa mengalikan dua bilangan prima besar sangatlah mudah, sementara memfaktorkan hasil kalinya kembali menjadi dua prima asal adalah masalah yang sangat sulit secara komputasi.

\subsection{Definisi dan Sifat Bilangan Prima}

Bilangan prima adalah bilangan bulat positif yang lebih besar dari 1 dan hanya memiliki dua pembagi positif: 1 dan dirinya sendiri. Bilangan bulat positif yang lebih besar dari 1 dan bukan prima disebut sebagai bilangan komposit.

Kriptografi menggunakan bilangan prima yang sangat besar, seringkali mencapai ribuan bit. Menghasilkan bilangan prima besar membutuhkan proses yang efisien karena kita tidak bisa menggunakan metode pembagian biasa (\textit{Trial Division}) yang akan memakan waktu jutaan tahun untuk angka sebesar itu.

\subsection{Algoritma Pengujian Primalitas}

Untuk menemukan bilangan prima besar, kita tidak mencari bilangan prima, tetapi memilih angka acak besar dan kemudian \textbf{mengujinya} apakah angka tersebut prima.

\begin{enumerate}
    \item \textbf{Sieve of Eratosthenes}: Metode klasik untuk menemukan semua bilangan prima dalam rentang tertentu. Namun, metode ini tidak efisien untuk angka yang sangat besar karena membutuhkan memori yang masif.
    \item \textbf{Pengujian Probabilistik (Miller-Rabin Test)}: Sebagian besar sistem kriptografi menggunakan pengujian probabilistik. Algoritma Miller-Rabin tidak memberikan kepastian 100\%, tetapi dapat memberikan jaminan bahwa suatu bilangan adalah "kemungkinan besar prima" (\textit{probably prime}) dengan tingkat kesalahan yang sangat kecil (misalnya $2^{-100}$). Jika suatu bilangan gagal dalam satu putaran uji Miller-Rabin, maka ia pasti komposit. Jika lolos, ia kemungkinan besar prima.
\end{enumerate}

\subsection{Pentingnya Bilangan Prima dalam Kriptografi}

Penggunaan bilangan prima bukan tanpa alasan. Sifat-sifatnya memberikan jaminan keamanan:
\begin{itemize}
    \item \textbf{Asimetri Komputasi}: Perbedaan kecepatan antara perkalian (mudah) dan faktorisasi (sulit) menciptakan "pintu satu arah" (\textit{trapdoor function}) yang menjadi dasar RSA.
    \item \textbf{Struktur Grup}: Bilangan prima $p$ memastikan bahwa himpunan $\{1, 2, \dots, p-1\}$ membentuk grup multiplikatif modulo $p$, yang menjamin bahwa setiap elemen memiliki invers modulo.
\end{itemize}

\begin{catatan}
Dalam implementasi nyata, bilangan prima yang digunakan untuk RSA harus dipilih secara acak dan sangat besar agar terhindar dari serangan faktorisasi menggunakan algoritma seperti \textit{General Number Field Sieve} (GNFS).
\end{catatan}
